% =============================================================
%  NLP2027 投稿論文 (和文)
%  貢献: PoR/ΔE による LLM 出力評価フレームワークと HA48 における予備検証
%
%  TODO[NLP2027-template]:
%    NLP2027 公式テンプレート公開後 (2026 年秋頃)、以下の
%    documentclass / usepackage を公式版に置き換えること。
%    現状は jsarticle ベースの portable な暫定構成。
% =============================================================

\documentclass[a4paper,11pt,twocolumn]{jsarticle}

\usepackage[utf8]{inputenc}
\usepackage{graphicx}
\usepackage{amsmath, amssymb, amsthm}
\usepackage{bm}
\usepackage{booktabs}
\usepackage{hyperref}
\usepackage{url}
\usepackage{xcolor}

% 査読中匿名化: 著者情報は最終調整時に置換
\title{LLM 応答の意味的整合性評価\\
       --- PoR/$\Delta$E 指標と HA48 における予備検証 ---}
\author{Anonymous Author \\ Anonymous Affiliation \\ \texttt{anonymous@example.com}}
\date{}

\begin{document}
\maketitle

% ===== Abstract =====
% TODO[(d)]: Claude が abstract 200 字 + contribution 1 文の 3 案を提示し、
%            User が最終文言を選択して以下に確定する。
\begin{abstract}
% [PLACEHOLDER]
% 想定構成:
%   - 問題設定 (LLM 応答の意味的整合性をどう測るか) : 1 文
%   - 提案 (PoR/ΔE/L_sem の枠組み)              : 1-2 文
%   - 検証 (HA48 における予備実験)               : 1 文
%   - 主要結果 (相関係数 / verdict 分布)         : 1 文
本稿では、大規模言語モデル (LLM) の応答に対する意味的整合性を
定量評価するための枠組みとして、PoR (Point of Resonance) 座標と
$\Delta$E 距離による指標群を提案する。
\end{abstract}

% ===== 1. はじめに (Introduction) =====
% 役割:
%   - 問題設定: なぜ LLM 応答の意味的整合性評価が必要か
%   - 既存手法の限界 (詳細は related_work へ、ここは概観のみ)
%   - 本研究の貢献を 3 点で箇条書き
%   - 論文構成の予告 (1 段落)
%
% 分量目安: 1 ページ (2 段組換算で約 800 字)
%
% 貢献候補 (locked-in: A 案):
%   PoR/ΔE による LLM 出力評価フレームワークと HA48 における予備検証

\section{はじめに}
\label{sec:introduction}

% TODO: 第一段落 - 問題設定
% LLM 応答の品質評価における既存指標 (BLEU, BERTScore, G-Eval 等) の
% 限界 (構造的整合性と命題カバレッジの分離不能性) を提起

% TODO: 第二段落 - 提案概観
% PoR (Point of Resonance) による 2 軸分離と ΔE 距離指標の導入

% TODO: 第三段落 - 提案概観 + 貢献
% locked-in (Variant 1): 以下の 1 文を第三段落の冒頭に置く
%   「本稿では、LLM 応答の意味的整合性を構造完全性 (S 軸) と命題カバレッジ
%    (C 軸) の 2 軸で分離評価する PoR/$\Delta$E フレームワークを提案し、
%    HA48 における予備検証を通じて人手評価との対応を確認する。」
%
% 主要貢献 (3 項目):
% \begin{enumerate}
%   \item PoR (S 軸 / C 軸) による 2 軸評価フレームワークの提案と、
%         理想応答からの距離指標 $\Delta$E の定式化
%   \item 意味損失関数 $L_{\mathrm{sem}}$ の 3 項分解 (演算子 / 未知性 /
%         因果構造) による判定根拠の解釈可能性
%   \item 48 件の人手アノテーション (HA48) を用いた予備検証による、
%         提案指標と人手評価との対応の実証
% \end{enumerate}

% ===== 2. 関連研究 (Related Work) =====
% 役割:
%   - LLM 応答評価の 3 系統を整理し、本研究の位置付けを明示
%   - 査読で必ず聞かれる「既存手法と何が違うか」への先回り回答
%
% 分量目安: 1 ページ (2 段組換算で約 800 字)
%
% 詳細サーベイ・全比較表は docs/related_metrics_comparison.md を参照
% (本論文の external reference として機能、camera-ready で URL 化)

\section{関連研究}
\label{sec:related_work}

LLM/NLG 応答の自動評価指標は、計算機構の観点から
(a) 表層類似度系、
(b) 埋め込み類似度系、
(c) LLM 駆動評価系
の 3 系統に大別できる。本節では各系統の代表的指標を概観し、本研究が
どの空白を埋めるかを示す。

\subsection{表層・埋め込み類似度指標}
\label{sec:rw:surface_embedding}

BLEU~\cite{papineni2002bleu} および ROUGE~\cite{lin2004rouge} は
n-gram overlap に基づく古典的指標であり、計算コストの低さと再現性から
広く採用されてきたが、構造的整合性と内容的妥当性を単一スコアに混合する
という限界を持つ。BERTScore~\cite{zhang2020bertscore} と
BLEURT~\cite{sellam2020bleurt} は事前学習言語モデルの埋め込み距離または
回帰モデルにより paraphrase 耐性を獲得したが、出力は依然として単一
スカラ (BLEURT) または P/R/F1 の限定的分解 (BERTScore) であり、応答の
どの側面 (構造 / 内容) で失点したかの診断情報を提供しない。

\subsection{LLM 駆動評価}
\label{sec:rw:llm_judge}

近年、評価そのものに LLM を用いる手法が急速に普及している。
G-Eval~\cite{liu2023geval} は GPT-4 に評価基準と Chain-of-Thought
を与え、確率重み付きスコアを生成する。Prometheus~\cite{kim2024prometheus,
kim2024prometheus2} は GPT-4 由来データで open-source LM を fine-tune
し、rubric 条件付きスコアと自然言語フィードバックを返す。
RAGAS~\cite{es2024ragas} は RAG パイプライン特化で 4 軸 (faithfulness,
answer relevancy, context precision/recall) に分解する。MT-Bench を
含む LLM-as-Judge~\cite{zheng2023judging} は GPT-4 を judge として
holistic スコアを生成する。これらは柔軟性が高い反面、(i) 評価軸の
分離が prompt / rubric 設計に完全依存し architectural な保証を持たない、
(ii) judge LLM のバージョン更新で評価が変動する再現性問題、という
共通の限界を抱える。LLM 評価論全体の俯瞰は Chang らのサーベイ
\cite{chang2024survey} に詳しい。

\subsection{本研究の位置付け}
\label{sec:rw:positioning}

上記いずれの系統も、応答の構造完全性 (どの程度、文として壊れていないか)
と命題カバレッジ (どの程度、問いの核心を拾っているか) を**計算機構の
段階で分離して評価する**枠組みを持たない。表層・埋め込み系は両者を単一
スコアに混合し、LLM 駆動系は分離を prompt 設計に委譲するため、軸の
排他性も網羅性も保証されない。本研究は次の 3 点でこの空白を埋める:
(i) 構造完全性 (S 軸) と命題カバレッジ (C 軸) を 2 次元座標 PoR として
architectural に分離する、(ii) 意味損失関数 $L_{\mathrm{sem}}$ を
3 項 (演算子・未知性・因果構造) に分解することで、判定根拠を再利用可能な
診断 primitive として提供する、(iii) judge LLM を計算経路に持たず、
事前学習文埋め込み (SBert) と cascade matcher のみで評価を完結することで
判定の再現性を担保する。

% ===== 3. 提案手法 (Method) =====
% locked-in 反映:
%   - PoR / ΔE / L_sem を主軸として展開 (Variant 1)
%   - Phase 8 advisory は §3.4 末尾で「補足機構」として浅めに、過大主張禁止
%
% 数式の出典:
%   - PoR / ΔE / quality_score: docs/formulas.md, ugh_calculator.py
%   - L_sem 3 項分解: docs/semantic_loss.md, semantic_loss.py
%   - verdict 閾値: docs/formulas.md (HA48 確定値)
%   - Phase 8 advisory: docs/phase_e_verdict_integration.md, mode_grv.py

\section{提案手法}
\label{sec:method}

本研究は、LLM 応答 $t$ の品質を、参照応答との 2 軸座標 PoR と距離指標
$\Delta E$ で表現し、さらに意味損失関数 $L_{\mathrm{sem}}$ の 3 項分解
により判定根拠を解釈可能にする。本節では各構成要素の定式化を順に示す。

\subsection{PoR: 構造完全性と命題カバレッジの 2 軸座標}
\label{sec:method:por}

応答 $t$ を、構造完全性 $S \in [0, 1]$ と命題カバレッジ
$C \in [0, 1]$ の 2 軸からなる座標
$\mathrm{PoR}(t) = (S, C)$ で表現する。

\paragraph{S 軸 (構造完全性).}
応答が構造的に破綻していないかを、検出層が抽出した
$k$ 項の f-flag $f_k \in [0, 1]$ の加重平均で定量化する:
\begin{equation}
\label{eq:s}
S = 1 - \frac{\sum_{k} w_k f_k}{\sum_{k} w_k},
\end{equation}
ここで $w = \{w_{f_1}, w_{f_2}, w_{f_3}, w_{f_4}\}
= \{5, 25, 5, 5\}$ である。$f_2$ (用語捏造) に最大重み 25 を配置するのは、
応答の信頼性を最も損なう失敗モードが命題の捏造であるとの校正結果に
基づく。$f_4$ が未計算の場合は当該重みを除外し $\sum w_k = 35$ で
正規化する。

\paragraph{C 軸 (命題カバレッジ).}
問いの核心命題集合 $\mathcal{P}$ に対する応答の被覆率を:
\begin{equation}
\label{eq:c}
C = \frac{\mathrm{hits}(t, \mathcal{P})}{|\mathcal{P}|},
\end{equation}
で定義する。$\mathrm{hits}(\cdot)$ は cascade matcher
(Tier 1: tfidf bigram, Tier 2: SBert 文埋め込み類似度) により決定される。
$|\mathcal{P}| = 0$ (命題未提供) の場合は $C$ を未定義 (degraded) とする。

\subsection{$\Delta E$: 理想応答からの距離}
\label{sec:method:delta_e}

理想応答 $(S, C) = (1, 1)$ からの加重二乗和距離として
$\Delta E \in [0, 1]$ を:
\begin{equation}
\label{eq:delta_e}
\Delta E = \frac{w_S (1 - S)^2 + w_C (1 - C)^2}{w_S + w_C},
\end{equation}
で定義する。$w_S = 2,\ w_C = 1$ は HA48 校正で人手評価との相関を
最大化する重みである。さらに人手 5 段階評価との対応のため、
パラメータフリーな線形変換:
\begin{equation}
\label{eq:quality_score}
\mathrm{quality\_score} = 5 - 4 \cdot \Delta E \in [1, 5],
\end{equation}
を quality\_score として併記する ($\Delta E = 0 \Rightarrow 5$,
$\Delta E = 1 \Rightarrow 1$)。

\subsection{$L_{\mathrm{sem}}$: 意味損失の 3 項分解}
\label{sec:method:l_sem}

$\Delta E$ は単一スカラに収束するため、判定がどの誤差項に駆動された
かを直接読み取れない。これを補うため、意味損失関数:
\begin{equation}
\label{eq:l_sem}
L_{\mathrm{sem}} = \alpha L_P + \eta L_F + \delta L_G + \mathrm{(appendix\ 4\ \mbox{項})},
\end{equation}
を導入する。本研究では HA48 校正で単独 $\rho$ が有意であった主構成
3 項を中心に議論する (appendix 4 項 $L_Q, L_R, L_A, L_X$ は理論的整合性
のため保持し、本論文の主評価からは除外する)。

\paragraph{$L_P$ — 命題損失.} 命題の \emph{欠落} を測る:
\begin{equation}
L_P = 1 - C,
\end{equation}
これは C 軸の補数であり、命題未提供時は未定義となる。

\paragraph{$L_F$ — 用語捏造損失.} 偽命題の \emph{混入} を測る:
\begin{equation}
L_F = f_2 \in \{0,\, 0.5,\, 1\},
\end{equation}
$L_P$ が命題の欠落を捉えるのに対し、$L_F$ は文脈 $c$ の参照集合
$\mathcal{R}_s$ に存在しない用語・概念の混入を捉え、$L_P$ と直交する
独立の誤差軸を成す。

\paragraph{$L_G$ — 因果構造損失.} 応答内の語彙重力分布 (grv) を
3 項式 (drift / dispersion / collapse) で集約した値:
\begin{equation}
L_G = \mathrm{grv}(t) \in [0, 1],
\end{equation}
を用いる。長文応答での論点拡散や、複数論点の単一塊への collapse を
検出する。grv の詳細な定式化 (drift / dispersion / collapse の 3 項式)
は紙幅の都合により本稿では割愛する。

\paragraph{重み校正.}
HA48 ($n = 48$) 現行パイプライン snapshot で、3 項主構成 $L_P + L_F + L_G$
は人手評価との Spearman 相関 $\rho = -0.548$ ($p < 0.001$) を達成し、
現行 $\Delta E$ 単独 ($\rho = -0.482$, $p = 0.0005$) を上回る予測力を
示した。運用重みは $\alpha = 0.27,\ \eta = 0.21,\ \delta = 0.35$ で、
full-sample 最適値から LOO-CV で補正した保守的な値である。

\subsection{verdict 判定と補助的 advisory 機構}
\label{sec:method:verdict}

$\Delta E$ の閾値により、応答の判定 verdict を 4 値で離散化する:
\begin{equation}
\mathrm{verdict}(t) =
\begin{cases}
\mathrm{accept}     & \text{if } C \neq \mathrm{None} \land \Delta E \leq 0.10 \\
\mathrm{rewrite}    & \text{if } C \neq \mathrm{None} \land 0.10 < \Delta E \leq 0.25 \\
\mathrm{regenerate} & \text{if } C \neq \mathrm{None} \land \Delta E > 0.25 \\
\mathrm{degraded}   & \text{if } C = \mathrm{None} \lor \Delta E = \mathrm{None}
\end{cases}
\end{equation}
閾値 $0.10 / 0.25$ は HA48 校正で確定済みであり、本論文では再校正を
行わず固定値として用いる。

\paragraph{補助的 advisory 機構 (簡略).}
応答モード (definitional, comparative, evaluative ほか) に条件付けた
重力解釈ベクトル (mode\_conditioned\_grv) を補助信号として導入し、
$\mathrm{verdict} = \mathrm{accept}$ かつ collapse\_risk が高く
anchor\_alignment が低い場合に限り、advisory として
$\mathrm{rewrite}$ への降格を提示する。これは主判定 verdict を
書き換えず、補助的な flag として併走する。
HA48+accept40 の拡充データ ($n = 63$) で校正した閾値は
$\tau_{\mathrm{collapse}} = 0.28,\ \tau_{\mathrm{anchor}} = 0.80$ で、
拡充検証データ上で $\rho = 0.523$ (主判定 $\rho = 0.441$) を達成した。
本機構の本格評価とモード別 ablation は本論文の scope 外であり、
\S\ref{sec:conclusion} の Future Work で議論する。

% ===== 4. 実験 (Experiments) =====
% 役割:
%   - HA48 アノテーションデータでの予備検証
%   - ベースライン比較 (G-Eval / BERTScore 最低 1 本)
%   - 数値の正直さ: 公式論文と docs/validation.md の同期を保つ
%
% 分量目安: 2 ページ
%
% TODO: 並行 backlog 完了後に最終数値を埋める
%   - HA48 → HA100 拡充版での再評価
%   - 第二 annotator + IAA (Cohen's κ)
%   - ベースライン (G-Eval) 比較

\section{実験}
\label{sec:experiments}

% ----- 4.1 設定 -----
\subsection{実験設定}
\label{sec:exp:setup}

% TODO: 4.1.1 データセット
% HA48: 48 件の人手アノテーション付き LLM 応答
% (拡充後は HA100 に差し替え)
\subsubsection{データセット}
% [PLACEHOLDER]
% data/human_annotation_48/annotation_48_merged.csv

% TODO: 4.1.2 比較手法 (ベースライン)
\subsubsection{比較手法}
% [PLACEHOLDER]
% - BERTScore
% - G-Eval (LLM-as-judge)
% - 提案手法 (PoR/ΔE)

% TODO: 4.1.3 評価指標
\subsubsection{評価指標}
% [PLACEHOLDER]
% - 人手評価との相関 (Spearman / Pearson)
% - verdict 分布と人手 verdict の一致率
% - Cohen's κ (annotator 間)

% ----- 4.2 結果 -----
\subsection{結果}
\label{sec:exp:results}

% TODO: 表 1 - 主要結果 (相関係数)
% docs/validation.md の current snapshot から転記

% TODO: 表 2 - verdict 分布の混同行列

% TODO: 図 1 - PoR 座標上での応答プロット

% ----- 4.3 考察 -----
\subsection{結果の解釈}
\label{sec:exp:discussion}

% TODO: L_sem 分解の効用 (どの損失項が verdict に効くか)
% TODO: 既存指標との差分 (S/C 分離の意義)

% ===== 5. 議論と限界 (Discussion & Limitations) =====
% 役割:
%   - Self-Audit Principle の位置付け (locked-in: ここに置く、名称は弱める)
%   - 本研究の限界を率直に列挙 → 査読者の懸念を先回りで回収
%
% 分量目安: 1 ページ
%
% locked-in 反映:
%   - Self-Audit Principle → 「出力の意味的誠実性に対する自己整合性評価」
%     という neutral な表現で言及
%   - 「無意識的重力仮説 (UGH)」は理論名として 1 段落のみ言及

\section{議論}
\label{sec:discussion}

% TODO: 5.1 自己整合性評価への適用可能性
% Self-Audit Principle を「出力の意味的誠実性に対する自己整合性評価」
% として再フレーミング。LLM 出力の自己診断ループとしての応用を議論。
\subsection{自己整合性評価への適用可能性}
% [PLACEHOLDER]

% TODO: 5.2 理論的位置付け
% 本指標群は著者らの提唱する「無意識的重力仮説 (UGH)」の枠組みに
% 位置付けられる。詳細な理論的動機は別稿に譲る。
\subsection{理論的位置付け}
% [PLACEHOLDER - 1 段落のみ]

% ----- 限界 -----
\section{限界}
\label{sec:limitations}

% TODO: 5.3 限界 (率直に列挙)
\subsection{本研究の限界}
% [PLACEHOLDER]
% - サンプル規模: HA48 (拡充後 HA100) は予備検証規模
% - annotator 数: 単一 annotator (拡充計画で第二 annotator + IAA を追加予定)
% - 単一ドメイン: AI 応答監査タスクに限定、汎化性は今後の検証課題
% - ベースライン: G-Eval / BERTScore のみ、より広範な比較は今後
% - 閾値固定性: HA48 で校正、拡充版での再校正は行わない (再現性確認のため)

% ===== 6. おわりに (Conclusion) =====
% 役割:
%   - 貢献の再掲 (Introduction の貢献文を圧縮)
%   - 今後の展望 (拡張型戦略への自然な接続)
%
% 分量目安: 0.5 ページ

\section{おわりに}
\label{sec:conclusion}

本稿では、LLM 応答の意味的整合性を構造完全性 (S 軸) と命題カバレッジ
(C 軸) の 2 軸で分離評価する PoR/$\Delta E$ フレームワークを提案し、
意味損失関数 $L_{\mathrm{sem}}$ の 3 項分解 (命題損失 $L_P$、用語捏造
損失 $L_F$、因果構造損失 $L_G$) により判定根拠の解釈可能性を担保した。
48 件の人手アノテーションデータ HA48 における予備検証では、$\Delta E$
が人手品質スコアとの間で Spearman $\rho = +0.482$ ($p < 0.001$) の
有意な相関を示し、BLEU・BERTScore・SBert cosine の 3 ベースラインを
点推定で上回った。また $L_{\mathrm{sem}}$ は $\rho = -0.548$ と
$\Delta E$ を上回る予測力を示し、主構成 3 項がそれぞれ独立した誤差軸
として寄与することを確認した。verdict 4 値判定の単調性も成立し、
閾値校正が人手判断と整合的に機能することを支持する結果を得た。

\paragraph{今後の展望.}
本論文の限界 (\S\ref{sec:limitations}) に対応して、以下 3 点を今後の
課題とする。(i) アノテーションを $n = 100$ 規模に拡充し、第二
アノテータを導入して Cohen's $\kappa$ により inter-annotator
agreement を定量化する。これによりベースラインとの相関差の統計的
有意性を $\alpha = 0.05$ 水準で確立する。(ii) $\Delta E$ の verdict
を補正する mode\_conditioned\_grv による advisory 機構
($n = 63$ 校正で $\rho = +0.523$) の本格的な ablation と、モード別
評価を別途実施する。(iii) 要約品質評価および対話応答評価という他
ドメインへの適用可能性を検証し、本フレームワークの汎化性を示す。

% ===== 謝辞 (Acknowledgments) =====
% TODO[最終調整]: 査読中は省略、camera-ready で追加

% ===== Data Availability =====
% TODO[最終調整]: locked-in 方針
%   「コードおよび HA48 アノテーションデータは、採択後に GitHub 上で
%    MIT ライセンスで公開予定である。査読中のアクセスが必要な場合は
%    program chair を経由して著者に連絡されたい。」


\bibliographystyle{plain}
\bibliography{references}

\end{document}
